\section{Discussion}
\label{sec:disc}
\paragraph*{What the experiments establish}
For stage V1, three statements are well supported. (1) Search on top of System~1 improves mean move quality modestly (\dAcplAbs\cp, about \fixedPct\% of System~1's blunders repaired) at \costRatio$\times$ the cost, through a small share of positions with large corrections. (2) System~1's confidence is weakly informative about whether its move is good (AUROC \aucConf) and does not identify where search helps: the confidence gate is indistinguishable from random gating at four rates, and two further predictions (\Hfour{}, \Hfive{}) are contradicted. (3) The V1 hybrid is not shown to be stronger than System~1 alone. The data do not show that the gate is useless, only that any advantage is below about \maxAdvHi\cp, and they do not show that selective search cannot work: the oracle indicates a perfect selector would exceed the full-search benefit while searching a few percent of positions, so what failed was one particular selector.

\paragraph*{Predictive uncertainty is not decision-criticality}
The exploratory evidence suggests a mechanism, offered as a hypothesis. Search helps rarely (\fracGainPos\% of positions improve) and in a heavy-tailed way: a few percent of positions, mostly forcing, endgame or tactically sharp, account for the net gain. There the model's ranking is sharply peaked, so confidence is \emph{high}, while positions with a flat ranking have little to gain. A confidence gate therefore sends compute where search changes the move often but with little net benefit. This is the difference between $P(\text{wrong})$, which confidence estimates, and $P(\text{search changes the move})\times\text{gain}$, which a router needs. Better-calibrated uncertainty~\cite{lakshminarayanan2017,gal2016} would improve the first without necessarily helping the second.

\paragraph*{What the follow-up stages change}
The follow-up alters the practical picture without overturning the V1 verdict. First, the largest improvement in move quality came from a component that is not a gate: a material-only blunder filter cut ACPL by 38\cp\ on fresh positions, far more than V1's 4.8\cp\ search effect (a different test set and agent, so not a controlled comparison). Its playing-strength gain is about 140 Elo alone but is not separated from the rules bundled with it, and the filter-free agent reached a similar rating. Second, with the repaired agent search added a further 8--11\cp\ over the filter alone, consistent with, but not proof of, the V1 conjecture that a root-prior temperature of 0.01 left search inert; the temperature, quiescence, tactical priors and value network changed together, and the old value network performed no differently from the new one. Third, a learned gate with access to System~1's win-chance margins, tactical flags and material was ahead of random gating at three rates: not significantly under the pre-registered test and significantly under a post hoc test with the same point estimates but a narrower interval. The effect, about 2.5\cp\ at a 25\% rate, is small beside the oracle headroom of about 20\cp. Fourth, System~1 matters once the filter is present (+212 Elo over a classical policy). Move quality, playing strength, compute efficiency and gating effectiveness were measured separately, and only the first two improved demonstrably.

\paragraph*{Compute, puzzles and confounders}
The V1 trade-off is nearly linear: the random gate's curve is the line between \ag{laya_only} and \ag{always_s2}, and the confidence gate lies on it. System~1 is the larger fixed cost (\msLayaPart\,ms of \msSTwo\,ms), so even a perfect gate removes only the search share. The small puzzle effect is consistent with search being anchored to System~1, and the tournament adds little power because the three System~1-based agents mostly draw each other. Confounders include the 0.01 prior temperature, the small value network, fixed-node judge noise in CPL, and the confound between position difficulty, confidence and game phase (openings have low confidence and few blunders, which contributes to the reversal of \Hfour{}).
